% !TeX root = ../papers/nc-p3.tex
\section{対角交換子零の全係数非零窓と二つの反対鉛筆}
\label{sec:marked-klein-zero-bracket-torus}
\begin{theorem}\label{thm:marked-klein-zero-bracket-torus}
対角差零の四元群交換子枝で二十四係数が非零とする。全ての窓は二符号 $u,h\in\{1,-1\}$ と
$(r_1,r_2,r_3)=(1,x,y)$, $xy\ne0$ による下記の四表示で覆われる。
$h=u$ の二表示は全点で標準Koszul左右AS正則、最小周期1である。
$h=-u$ の二表示では正則性は $xy+x+y\ne0$ と同値で、正則点の周期は2を割る。
従って前二節と合わせ、これらの四元群次数を持つ印付き交換子表示の全台・全パラメータでAS性が決まる。
\end{theorem}
\begin{proof}[二符号への縮約]
$Q_i^- =1,Q_i^+=u$ と取り、対角差零より $P_i^-=r_i,P_i^+=ur_i$ と書く。
共通倍率で $r_1=1$ とする。全係数非零なので定理\ref{thm:marked-klein-mixed-torus}の
非対角の計算を $v_i=u$ で使える（この段では対角差で割らない）。
二つの対角外三項積の一致は $u^6=1$ を与え、軸の対角比で同じ三乗根を掛けて $u=\pm1$ にできる。
$G_i^+=c_i,G_i^-=hc_j$（以下 $i,j,k$ は巡回）で
$c_j/c_k=r_i/r_k$ となるので共通倍率を選び $c_i=1/r_j$ にする。
$F_i^+=-r_kr_i,F_i^-=-r_ir_j/h$ となり、$G_i^+F_i^-=G_i^-F_i^+$ は $h^2=1$ である。
逆にこれらの係数は全交換子式を満たす。
\end{proof}

\begin{proof}[$h=-u$ の符号付き平方判定]
反復対の純成分は $w_{jkjk}=w_{kjkj}=u(r_j+r_k)$ で、対向星の対は1である。
まず $\prod(r_j+r_k)\ne0$ とする。定理\ref{thm:marked-klein-mixed-torus}と同じ
三十六行のWalsh規格化の整数左核積が全て1となり、同じ八行による余因子の比は
\[
       \frac{N_{11}N_{22}}{N_{12}N_{21}}-1
            =\frac{(x+y)(y+1)}{y^2}-1
            =\frac{xy+x+y}{y^2}.
\]
三行三列は非零で行和・列和が零なので、これが零なら階数1、非零なら階数2である。
各 $r_j+r_k=0$ のときは $r_1r_2+r_2r_3+r_3r_1=-r_j^2\ne0$。
前節までと同じく欠けた色ごとに四位置の $0,i$ を $z$ 倍して $z^2$ で割り、
極限で色対 $(i,i)$ を全て消す。残った整数左核積を特殊化でき、非対角の全成分が非零なので
三色の階数2ブロックと階数2の符号付き平方行列が残る。Walshの定理と良還元によって元の窓も正則である。
二位置移動の周期性は分母が消える点も元の係数で確認できる。
実際移動後の $Q_i^+,Q_i^-$ は $r_i,ur_i$、二対の局所積と $Q_i^+/Q_i^-$ の積は変わらない。
$G$ 側も比の積 $h$ を保つ（$F_i$ は $G_i$ の倍数、転置で比は逆数）。
上の導出で $u,h$ はこの二つの積、$x,y$ は $Q_i^-P_i^+$ の射影比で決まるので元の標準表示に戻る。
\end{proof}

\begin{proof}[$u=h=1$ の両反対核と全評価]
位置2を $\operatorname{diag}(1,-1/(xy),-1/(xy),-1/(xy))$、位置3を
$\operatorname{diag}(1,1,1/x,1/y)$ で変更すると両対称テンソルになる。
反対切断の対称二乗の三つの非自明色の二次行列式はそれぞれ $1/(xy)^2,1/(xy^2),1/(x^2y)$ である。
零色の四次行列は階数2で、その右核と左核の二基底は下の行列の列になる。
\[
 \begin{pmatrix}
 -1+1/x&-1+1/y\\ -1/x&-1/y\\1&0\\0&1
 \end{pmatrix},\qquad
 \begin{pmatrix}
 (1-x)/(xy)&(1-y)/(xy)\\-1&-1\\1&0\\0&1
 \end{pmatrix}.
\]
従って反対階数は全点で8であり、両側とも対角二次式の鉛筆である。
四行のどの二行も独立になる条件は $1,x,y$ が相異なることである。
このとき四つの特異対角二次式は相異なる単根に対応し基底曲線は滑らかな完全交叉である。
実際特異点なら勾配の従属する鉛筆元の核に入るが、階数3の対角元の核は座標軸であり
他の鉛筆元はそこで非零である。重なる場合は二行を同時に消す元が取れて階数2以下の二次式を含む。

第一辺の純関係が存在しないことを示す。非零な純関係の星・星成分は零なので少なくとも一端の星成分が零になる。
星行が零だと星行を含む三関係の係数が零になり、残る三関係の一次結合は列を $a,Q_i a$ とした行列になる。
階数1なら $ka$ は三行列の共通固有直線である。逆端では $G_i$ またはその転置が同じ役割を持つ。
いずれの組でも各第 $i$ 行列の平方は $i$ 軸以外への座標射影の非零倍で、
全軸の間は非零の行列成分で連結するので共通不変直線はない。
滑らかな場合は定理\ref{thm:smooth-apolar-complete}で正則性が従う。

重なる場合のために全評価の階数下界を示す。点 $[t:a]$ での評価の核を $(s,v)$ とすると
\[
       (tI+G_a^t)v=0,\qquad s a+Q_v^t a=0 .
\]
$a=0$ なら核は一次元、$a\ne0$ なら $v$ から $s$ が高々一通りに決まる。
最初の行列が階数2以上ならこれでよい。$t=0$ で $a\ne0$ のとき、
各非零な成分 $a_i$ に対応する二行二列主小行列式が非零なので階数1以下にならない。
$t\ne0$ で階数1なら、対角三成分が全て非零なので全非対角成分も非零、従って $a_1a_2a_3\ne0$。
ところが $v\mapsto Q_v^t a$ の行列式は $2a_1a_2a_3$ なので、これが $ka$ に入る $v$ は高々一次元である。
他端でも同様で、対応する行列式は $2a_1a_2a_3/(xy)$ である。全評価が階数3以上と分かった。
補題\ref{lem:bisymmetric-reducible-quick-test}で重複の場合も全次数正則である。
\end{proof}

\begin{proof}[$u=h=-1$ と両表示の周期]
四元群の二文字積に双指標 $\epsilon(a,b)=(-1)^{a_1 b_0}$ を掛けるコサイクル捻りを行う。
各多重次数のバー複体で語 $(a_1,\ldots,a_m)$ を $\prod_{i<j}\epsilon(a_i,a_j)$ 倍することで
積の微分と捻った微分が一致する。Hilbert級数とKoszul性を保ち、二次双対の
Frobenius対合も各相補次数で非零倍率を掛けるだけなので左右AS性も保つ。
四テンソルでも同じ積を掛け、位置2,3のラベル $a$ を $\epsilon(a,a)$ 倍して星の単位成分を戻す。
得られる $Q_i^+/Q_i^-$ の積と $G_i^+/G_i^-$ の積はともに1、$Q_i^-P_i^+$ は順に $1,x,y$ である。
従って上の規格化により $u=h=1$ の解そのものになり、逆の捻りで正則性を得る。

最後に両符号の周期1をテンソルの軸の混同なしに確認する。位置 $p$ の三つの非星文字に
倍率を選び、全四位置の六関係を同じ六式
\[
       t x_i+A_i x_jx_k+B_i x_kx_j=0,\qquad
       x_i t+u B_i x_jx_k+u A_i x_kx_j=0
\]
に揃えられる。これは十二倍率と六係数を未知とする四十八の指標式であり、
指数行は各非星積の二倍率、星積の非星文字の逆倍率、該当する $A_i,B_i$ の逆数からなる。
全ての右辺は非零Laurent単項式であり、Smith零行の全指標が1であるため全点で解がある。
\path{computations/verify_klein_zero_bracket_torus.py} はこの二符号の指標式、Walshの整数左核、
反対鉛筆、評価の二行列式、コサイクルの符号を検算する。
\end{proof}

なお前節の疎な正則表示も二位置移動で戻る。実際 $Q'_i=P_i^t,P'_i=Q_i^t,G'_i=F_i^t,F'_i=G_i^t$
を前節の表示で取り、第二枠に対角倍率 $(-b/x,1,1/x)$ を掛け、第一枠は単位とする。
$Q_2,Q_3$ が元の1に、$G_2,G_3$ の正成分も1に、負成分も $b,x/b$ に戻る。
従って四元群交換子表示が一つでも存在する標準窓は、対角差と台によらず周期が2を割る。
印を忘れたときの最小周期の残る二択は、同位置の単純印を全て比較することで決まる。
実際星点の局所一次式は六独立であり単純印は有限個、残る枠の同一視は定理\ref{thm:marked-point-commutators}の
二つの可逆三行列と一倍率の式が必要十分である。対角差零では枠を単項行列に限る必要はない。
